\documentclass{beamer}

\usepackage[T2A]{fontenc}
\usepackage[utf8]{inputenc}
\usepackage[english,russian]{babel}

\usepackage{tikz}

% \usepackage{mathpazo}      % шрифт Palatino + математика
\usefonttheme{serif}       % beamer использует шрифт с засечками

\usetheme{CambridgeUS}
\usecolortheme{wolverine}

\setbeamertemplate{caption}[numbered]

% \usepackage{pgfpages}
% \setbeameroption{show notes on second screen}

\title[Математическое моделирование]{Математическая модель}
\author[Гркикян М.Э.]{Гркикян Мисак Эдикович}
\date{09.10.2026}

% \setbeamertemplate{title page}{
% 	\centering
% 	\vskip1em
% 	{\usebeamerfont{title}\usebeamercolor[fg]{title}\inserttitle\par}
% 	\vskip1em
% 	{\usebeamerfont{author}\insertauthor\par}
% }

\begin{document}
    \begin{frame}
        \titlepage
    \end{frame}

    \begin{frame}{Модель расхода воды в душе}
        \textbf{Этап 1. Формулировка предмодели.}
        \begin{itemize}
            \item Объект - процесс расхода воды в душе.
            \item Цель - определить, сколько воды утекает за одно принятие душа.
            \item Предположения: напор постоянный, утечек нет.
        \end{itemize}

        \medskip
        \textbf{Этап 2. Завершение идеализации.}
        \begin{itemize}
            \item Пренебрегаем испарением, нагревом от среды, трением в трубах.
            \item Математически:
                $V_{\text{исп}} = 0$,\;
                $p = \text{const}$,\;
                $T_{\text{смеси}} = \text{const}$.
        \end{itemize}
    \end{frame}

    \begin{frame}{Модель расхода воды в душе}
        \textbf{Этап 3. Выбор закона.}
        \begin{itemize}
            \item Объект подчиняется \textbf{закону сохранения массы}:
        \end{itemize}
        \[
            V(0) + V_{\text{прих}} = V(t) + V_{\text{расх}}.
        \]
        \begin{itemize}
            \item Дополнительно: расход пропорционален времени,
        \end{itemize}
        \[
            V_{\text{расх}} = q_{\text{расх}} \cdot t.
        \]

        \medskip
        \textbf{Этап 4. <<Оснащение>> модели.}
        \begin{itemize}
            \item $V(0) = 0$ л (бак пуст), $q_{\text{прих}} = 12$ л/мин (сколько воды приходит за минуту),
                $q_{\text{расх}} = 12$ л/мин, $t = 10$ мин.
            \item Цель - найти $V_{\text{расх}}$.
        \end{itemize}
    \end{frame}

    \begin{frame}{Модель расхода воды в душе}
        \textbf{Этап 5. Исследование модели.}
        \begin{itemize}
            \item Вода идёт напрямую, бак не участвует:
                $V(0) = 0$, $V(t) = 0$, значит
                $V_{\text{прих}} = V_{\text{расх}}$.
        \end{itemize}
        \[
            V_{\text{расх}} = q_{\text{расх}} \cdot t = 12 \cdot 10 = 120\ \text{л}.
        \]
        \begin{itemize}
            \item Ответ: за 10 минут душа утекает 120 литров.
        \end{itemize}

        \medskip
        \textbf{Этап 6. Проверка адекватности.}
        \begin{itemize}
            \item Типичный душ - 10-15 л/мин, значит 100-150 л за 10 минут,
                результат правдоподобен.
            \item Не противоречит закону сохранения массы.
        \end{itemize}
    \end{frame}

\end{document}